% !TeX root = ../../../main.tex
\subsection{Survival Modeling with Language Model Embeddings}
\label{ch07:apd:encoder-survival}


% !TeX root = ../../../main.tex
\begin{table}[htbp]
\caption{Hyperparameter grid search ranges for our survival models.}\ChSevenFloatLabel{ch07:tab:hyperparams}
\ChSevenTableBodyStart
\begin{ChSevenFloatBody}

\centering
\begin{adjustbox}{max width=\columnwidth}
\begin{tabular}{ll}
\toprule
\textbf{Model} & \textbf{Hyperparameters} \\
\midrule
RSF & \texttt{n\_estimators} $\in$ \{20,50,100,200,500,1000\} \\
    & \texttt{min\_samples\_split} $\in$ \{2,10,20,50,100\} \\
    & \texttt{min\_samples\_leaf} $\in$ \{1,5,10,20\} \\
DeepSurv/DeepHit & \texttt{num\_nodes} $\in$ \{32,64,128,256,512,768,1024,1200,\\
& 1536,1800,2048,3072,4096\} \\
    & \texttt{dropout} $\in$ \{0.1,0.2,0.4,0.5\} \\
\bottomrule
\end{tabular}
\end{adjustbox}



\end{ChSevenFloatBody}
\end{table}

The embedding and survival-learner specifications are given in Section~\ref{ch07:sec:survival-modeling}.

Patients who do not die are right-censored at the timestamp of their last recorded event.
Patients whose death or last observed follow-up precedes the observation cutoff (0, 24, or 168 hours) are excluded.
Survival time is measured from the observation cutoff.

Hyperparameter tuning for each survival model is performed using a validation split from the training set. 
For each combination of language model and survival model, we conduct grid search over the respective hyperparameter space, selecting the configuration that maximizes the time-dependent concordance index \citep{ch07:antolini2005time} on the validation set. 
We use five-fold cross-validation on training data to select the best-performing survival learner for each annotation model.
Final evaluation is conducted on held-out test sets (\textit{T2S2\_test}, \textit{sepsis-10}, \textit{sepsis-100}), and results are aggregated across survival metrics to assess model performance.

This framework allows for a systematic evaluation of how different pre-trained LLMs, and their associated embedding strategies, contribute to downstream survival prediction tasks.
\FloatBarrier
